% !TeX root = ../nc-quadric-avatar.tex
% 2026-11-22: 共有根 A1 の最小解消に残る相対 Pic^0 障害
\section{共有根の最小解消と相対 Picard の奇偶障害}
\label{sec:qa-nodal-resolution-parity}

定理~\ref{thm:qa-noncubic-rootpair-torsionfree-spectral-module}
の相対 spectral module は共有根で非局所自由となる．
この非自由性は，その点を最小解消して線束にすれば
自動的に通常の相対 $\Pic^0$ に戻る，という意味ではない．
以下は元の一径数基底を変えない最小解消での具体的な障害である．

\begin{proposition}[最小解消後にも残る奇数次数]
\label{prop:qa-nodal-resolution-odd-exceptional-degree}
$k$ は標数零の代数閉体とし，
$A=k[[a,b,y]]/(ab-y^2)$，
$s=a-b$ とおく．
符号商
\[
 A=k[[c,d]]^{\{\pm1\}},\qquad
 a=c^2,\quad b=d^2,\quad y=cd
\]
の反不変部分
$M=k[[c,d]]^- =Ac+Ad$
を考える．$\pi:\widetilde X\to\Spec A$ を
$A_1$ 特異点の最小解消，$E\simeq\PP^1$ を
例外曲線とする．
\[
 L=(\pi^*M)/(\text{torsion})
\]
は $\widetilde X$ 上可逆で，その例外曲線への制限は
\[
                    L|_E\simeq\OO_{\PP^1}(1).
\]
特に任意の $n\in\mathbb Z$ に対し
$\deg_E L(nE)=1-2n$ は奇数であり，
相対 $\Pic^0$ に属する多重次数零の線束にはならない．

さらに，$s$ を基底パラメータとする適切な大域曲線族で
特殊曲線が既約一節点であれば，最小解消後の特殊
fiber は二つの成分 $D'$ と $E$ を持ち，
$D'\cdot E=2$ である．
同じ generic fiber 上で $M$ と一致する
任意の可逆延長は，$L$ から垂直 Cartier 因子で
tensor したものであるから，
$E$ 上の次数の奇偶は変わらない．
従って\emph{元の一径数基底上で最小解消するだけでは}，
非自明二捩れの generic fiber を
多重次数零の可逆線束の相対族として延長できない．
\end{proposition}
\begin{proof}
最小解消の標準二枚のアフィン chart は
\[
\begin{array}{c|ccc}
 &a&y&b\\ \hline
 U_c&u&ut&ut^2\\
 U_d&vw^2&vw&v
\end{array},\qquad
v=ut^2,\quad w=t^{-1}.
\]
例外曲線は $u=0$ と $v=0$ を貼り合わせた
$\PP^1_{[c:d]}$ で，自己交点は $E^2=-2$ である．
$M$ の生成元を形式的に $e_c=c,\ e_d=d$ と書く．
$A$ 上の関係 $yc=ad$ と $yd=bc$ を引き戻して
torsion を除くと，
$U_c$ では $e_d=t e_c$，
$U_d$ では $e_c=w e_d$ となる．
従って各 chart では一つの自由生成元を持ち，
重なりでの遷移関数は
$e_d=t e_c$ である．これは
$\OO_{\PP^1}(1)$ の標準遷移である．
$\OO(E)|_E=\OO_E(-2)$ だから
$\deg_E L(nE)=1-2n$ が従う．

$s=a-b$ を引き戻すと
$U_c$ では $s=u(1-t^2)$，
$U_d$ では $s=v(w^2-1)$ となる．
従って特殊 fiber には例外成分 $E$ が重複度一で現れ，
元の節点の二枝の strict transform が
$E$ の $t=1,-1$ の二点に横断的に交わる．
大域的に元の特殊曲線が既約一節点なら，
この二枝は同じ既約 strict transform $D'$ に属する．
ゆえに $D'\cdot E=2$ である．
二つの可逆延長が generic fiber 上同型なら，
正則曲面上のその比は垂直 Cartier 因子で表される．
その成分は $E$ と $D'$ の整数係数の和なので，
$E$ 上の次数の変化は
$nE^2+mD'\cdot E=-2n+2m$ で偶数である．
よって多重次数零への延長は存在しない．
\end{proof}

\begin{corollary}[相対 avatar の修正に必要な追加操作]
\label{cor:qa-nodal-resolution-requires-more}
共有根の noncubic spectral avatar を
通常の相対 $\Pic^0$ へ移すには，
元の基底を変えずに全空間の $A_1$ を最小解消するだけでは
不十分である．
基底変更，非可逆 rank-one torsion-free 層を許す
compactified Picard，または安定モデルと偏極を
同時に変更する構成のいずれかを別途検討する必要がある．
この命題は，それらの追加操作を施した場合にも
相対射が存在しないと主張するものではない．
\end{corollary}
\begin{proof}
命題~\ref{prop:qa-nodal-resolution-odd-exceptional-degree}
で可逆延長の例外成分上の次数が奇数になる．
相対 $\Pic^0$ の多重次数零の線束にはなり得ない．
\end{proof}


\subsection{Spectral 二捩れ乗法の欠損と正規化上の極限}
\label{subsec:qa-spectral-compactified-limit}

前の奇偶障害は，線束の極限が存在しないという意味ではなく，
\emph{可逆線束に限った極限が存在しない}ことを意味する．
相対 torsion-free sheaf 自身には自然な二次積があり，
その余核によって境界を scheme-theoretic に測定できる．
以下では定理~\ref{thm:qa-noncubic-rootpair-torsionfree-spectral-module}
の記号を使い，二形式は各々二つの相異なる根を持つとする．

\begin{proposition}[乗法欠損の長さと自己双対性]
\label{prop:qa-spectral-product-defect}
$B=\PP^1_k$，$f_1,f_2\in H^0(B,\OO_B(2))$ は
非零で各々 squarefree，$k$ は標数零の代数閉体とする．
$C=\{y^2=f_1f_2\}$ と
$\mathcal M=(r_*\OO_{\widetilde C})^-$ を取る．
反不変元の積から生じる
\[
 \mu:\mathcal M\otimes_{\OO_C}\mathcal M\longrightarrow\OO_C
\]
について，次が成り立つ．
\begin{enumerate}
\item $\operatorname{im}\mu=\mathcal I=(f_1,f_2,y)$ であり，
  $T=V_B(f_1,f_2)$ を $y=0$ によって $C$ に埋め込むと
  \[
     \coker\mu\simeq\OO_T,\qquad
     \operatorname{length}(\coker\mu)
       =\deg\gcd(f_1,f_2)\in\{0,1,2\}.
     \tag{Sq-defect}\label{eq:qa-spectral-square-defect}
  \]
  この余核は二根対の共通零点の
  \emph{scheme-theoretic} な構造層である．
\item 各幾何 fiber において，積の随伴射
  \[
   \mu^\flat:\mathcal M
       \longrightarrow
       \mathcal H om_C(\mathcal M,\OO_C)
  \]
  は同型である．つまり $\mathcal M$ は自己双対だが，
  $\mu$ が全射となるのは二根対が交わらない場合に限る．
\end{enumerate}
特に $\mu^\flat$ の同型と
可逆線束としての $\mathcal M^{\otimes2}\simeq\OO_C$
とは，節点では異なる主張である．
\end{proposition}
\begin{proof}
相対射 $r:\widetilde C\to C$ に関する
同時符号の反不変元は局所的に $c,d$ で生成され，
$c^2=f_1,\ d^2=f_2,\ cd=y$ である．
従って積の像は記載のイデアルで，
\[
 \OO_C/(f_1,f_2,y)
 \simeq\OO_B/(f_1,f_2)
\]
が従う．$f_i$ は各々 squarefree なので，
その共通零点は $0,1,2$ 個の被約点である．

自己双対性は，二形式が共通根を持たない開集合では
通常の非分岐二重 torsor の二次積から従う．
共通根では完備局所座標と単元の平方根を選ぶと，
\[
 A_0=k[[x,y]]/(y^2-x^2),\quad
 M_0=\langle c,d\mid yc=xd,\ yd=xc\rangle,
 \qquad c^2=d^2=x,\quad cd=y .
\]
ここで $u=c+d,\ v=c-d$ とすれば
\[
  M_0\simeq A_0/(y-x)\,u
          \ \oplus\ A_0/(y+x)\,v,
  \qquad u^2=2(x+y),\quad
  v^2=2(x-y),\quad uv=0.                 \tag{Node-M}
\]
環 $A_0$ において
\[
 \begin{split}
  \Hom_{A_0}(A_0/(y-x),A_0)&=(x+y),\\
  \Hom_{A_0}(A_0/(y+x),A_0)&=(x-y)
 \end{split}
\]
であり，積からの $u\mapsto 2(x+y)$ と
$v\mapsto2(x-y)$ は各直和因子で同型となる．
異なる枝間の Hom は零なので，$\mu^\flat$ は同型である．
各共通根は同じ局所型であり，主張が従う．
\end{proof}

\begin{theorem}[noncubic 三型の compactified spectral limit]
\label{thm:qa-spectral-three-fiber-limits}
前命題の仮定を置き，$\nu:\widehat C\to C$ を
節点での正規化とする．
\begin{enumerate}
\item 二根対が交わらない smooth 型では
  $\mathcal M$ は非自明な $\Pic^0(C)[2]$ の線束であり，
  $\End_C(\mathcal M)=k$ である．
\item ちょうど一根共有する既約 nodal 型では
  $\widehat C\simeq\PP^1$ であり，
  \[
       \mathcal M\simeq\nu_*\OO_{\PP^1}(-1),
       \qquad \End_C(\mathcal M)=k .
       \tag{Node-limit}\label{eq:qa-spectral-nodal-limit}
  \]
  これは次数零の rank-one torsion-free \emph{単純}層だが，
  $\Pic^0(C)\simeq\Gm$ の非自明二捩れ線束
  （節点の貼合せ比 $-1$）とは同型でない．
\item 二根対が一致する $I_2$ 型では
  $\widehat C=C_+\amalg C_-$，$C_\pm\simeq\PP^1$ であり，
  \[
       \mathcal M\simeq
       \nu_*(\OO_{C_+}(-1)\oplus\OO_{C_-}(-1)),
       \qquad \End_C(\mathcal M)\simeq k\oplus k .
       \tag{I2-limit}\label{eq:qa-spectral-i2-limit}
  \]
  従って $\Aut_C(\mathcal M)\simeq\Gm^2$ で，
  この層は\emph{単純でない}．
\end{enumerate}
三型とも $\chi(\mathcal M)=0$ である．
従って，前定理の全相対族を受けるモジュライを
\emph{単純な} rank-one torsion-free sheaf に限定すると，
$I_2$ 型の点が脱落する．
ただし適切な偏極に関する半安定層の粗モジュライを
排除するものではない．
\end{theorem}
\begin{proof}
(1) は定理
~\ref{thm:qa-noncubic-rootpair-torsionfree-spectral-module}
により従う．
各節点での式 (Node-M) は
正規化の二枝の直像の局所模型そのものである．
従って非局所自由な全節点で正規化すると，
$\mathcal M=\nu_*Q$ となる線束 $Q$ が存在する．

(2) の正規化は種数零の滑らかな連結曲線であり，
$Q\simeq\OO_{\PP^1}(d)$ と書ける．
既知の
$p_*\mathcal M\simeq\OO_{\PP^1}(-1)^{\oplus2}$
によって $\chi(\mathcal M)=0$ なので，
$\chi(Q)=d+1=0$，従って $d=-1$ である．
また
$\mathcal E nd_C(\nu_*Q)\simeq\nu_*\OO_{\widehat C}$
であるから，大域 End は $k$ となる．
$\Pic^0(C)$ の元は節点でも局所自由なので，
非局所自由な $\mathcal M$ とは同型にならない．

(3) では比例する二形式をスカラーで規格化して
$f_1=f_2=f$ とできる．曲線は $y=\pm f$ の二成分で，
正規化後の各成分は $B\simeq\PP^1$ に同型である．
$y$ の $c,d$ 上の作用は
$\left(\begin{smallmatrix}0&f\\f&0\end{smallmatrix}\right)$
なので，$c+d$ と $c-d$ がそれぞれ
$y=f$ と $y=-f$ に支持される固有生成元となる．
$p_*\mathcal M=\OO_B(-1)c\oplus\OO_B(-1)d$ から
各固有成分は $\OO_B(-1)$ である．
異なる既約成分間には非零な sheaf Hom はなく，
各成分の自己準同型環は $H^0(\OO_{\PP^1})=k$ だから
End は $k\oplus k$ となる．
最後の留保は，simple/fine と
semistable/coarse の compactification を
混同しないためである．
\end{proof}

\begin{remark}[相対 avatar に必要なモジュライの選択]
\label{rem:qa-spectral-compactification-choice}
既約 nodal 型では $\nu_*\OO_{\PP^1}(-1)$ が
単純な compactified Picard の境界点となる一方，
$I_2$ 型では automorphism が $\Gm^2$ になる．
rank-one torsion-free sheaf の compactification と
安定性条件の一般論は
\cite{AltmanKleiman1980,ChristPayneShen2023} を参照する．
本稿の固有な入力は式~\eqref{eq:qa-spectral-square-defect}と
二根対の三型に沿った explicit な極限表示である．
今後の正確な比較対象は，
相対曲線 $C/B$ と $\mathcal M$ のみならず，
\emph{自己双対な二次積 $\mu^\flat$ と
長さ $0,1,2$ の積欠損}も保持する群oidである．
この群oidと avatar AS 心の同値は未証明である．
\end{remark}
